%! TeX root = ./main.tex
% vim: set spelllang=en:
\section{Multiple reconciliation}

Some hosts carry several families of units.
To reconstruct the respective histories of these families, bilateral reconciliation methods could be applied to them individually.
This approach would, however, assume that the units under study evolve independently of one another, which is in particular not the case for some genes.
Genes are indeed organized linearly along the genome of their species and affected by genome-scale events.
These events modify the organization and the composition of the genome, acting simultaneously on several of its genes.

The evolution of linear structures, shaped by rearrangements and duplications, is mainly studied in the context of comparative genomics.
The \emph{genome phylogeny} problem consists in reconstructing the most plausible phylogenetic tree for a set of genomes, by associating each internal node with an ancestral genome.
This problem comes in a \emph{\enquote{large}} version, where the topology of the tree and the ancestral genomes are inferred simultaneously, and a \emph{\enquote{small}} version, in which the topology is fixed and only the ancestral genomes remain to be determined~\parencite[for a review of these problems, see][]{ElMabrouk2012AnalysisGeneOrder}.

To address these problems, \textcite{Sankoff2000DuplicationRearrangementReconciliation} proposed combining the methods developed to analyse evolution at the scale of gene families with those for the scale of genomes.
The information recovered by reconciliation, in particular orthology relationships, could indeed guide the study of the evolution of genomes containing several copies of the same gene family.
Conversely, integrating the genomic scale into reconciliation would make it possible to take into account the interdependence of the evolution of genes.
The research objective is thus the development of methods capable of reconstructing the joint history of several gene families~\parencite[see also][]{Chauve2013DuplicationRearrangementReconciliation}.

To do so, an intermediate level can be integrated into the bilateral reconciliation model, placed between the scale of genes and that of species, in order to model evolution at the genomic scale.
In what follows, we compare two existing approaches based on this idea according to their modelling capacity and their algorithmic complexity.

\subsection{Adjacency reconciliation}

The linear organization of a set of units can be represented by a set of adjacencies between pairs of units.
For example, the genome \(\{(a, b, c), (d, e)\}\) made up of two chromosomes and five genes corresponds to the three adjacencies \(\{\{a, b\}, \{b, c\}, \{d, e\}\}\).
This representation notably makes it easier to model rearrangement events, since it suffices to act on the adjacencies at the border of the segment moved by the event.
The DeCo approach (for \enquote{Detection of Coevolution}), introduced by \textcite{Berard2012EvolutionGeneNeighborhoods}, extends reconciliation to integrate the evolution of adjacencies between pairs of dependent units.

Suppose that a forest \(\gtrees\) of phylogenies of dependent units and the phylogeny \(S\) of their hosts are known.
An adjacency \(\{a, b\}\) between two internal nodes of \(\gtrees\) is an assertion stating that the two divergences occurred simultaneously in the same host and that the copy of \(a\) is an immediate neighbour of \(b\) in the structure of this host.

\begin{definition}[Adjacencies]\label{definition:adj}
    An \emph{adjacency} on a set \(L\) is a pair \(\{a, b\} \subseteq L\) such that \(a \neq b\).
    A map \(f : L \onto M\) \emph{respects} such an adjacency if \(f(a) = f(b)\).
    The set of all possible adjacencies respected by~\(f\) is denoted~\(\alladjs(f)\).
\end{definition}

Suppose that the set of adjacencies between the leaves of \(\gtrees\) is also known.
An adjacency reconciliation of \(\gtrees\) with \(S\) (see \cref{fig:reconciliation:adj-rec}) is a reconciliation of each tree of \(\gtrees\) with \(S\) together with a set of ancestral adjacencies leading to the adjacencies between the leaves.

\begin{definition}[Adjacency reconciliations]\label{definition:adjrec}
    Let \(\gtrees\) be a phylogenetic forest and \(S\) a phylogenetic tree.
    Let \(\llmap : \tleavs(\gtrees) \onto \tleavs(S)\) be a map and \(\lladjset \subseteq \alladjs(\llmap)\) a set of adjacencies.
    An \emph{adjacency reconciliation} of \(\gtrees\) with \(S\) is a pair \(\tup{\rec, \adjset}\),
    where \(\rec\) is a map extending~\(\llmap\) that is the union of individual reconciliations of each tree of \(\gtrees\) with \(S\),
    and \(\adjset\) is a set of ancestral adjacencies such that \(\lladjset \subseteq \adjset \subseteq \alladjs(\rec)\).
\end{definition}

\begin{figure}[htb]
    \capstart
    \centering
    %! TeX root = ../main.tex
% vim: set spelllang=fr:
\begin{tikzpicture}
    %% Arbre T_a
    \begin{scope}[color1]
        \graph[
            phylogenetic tree layout,
            sibling distance=.2cm,
            evolutionary unit length=0.44cm,
            nodes={font=\small},
        ]{
            ta l4s1 / [coordinate] -- {
                ta l3s2 / [coordinate] -- {
                    ta l2s2 / [coordinate] -- {
                        ta l1s6 / \(\explar{a}{3}{2}\),
                        ta l1s5 / \(\explar{a}{2}{2}\),
                    },
                    ta l1s4 / \(\explar{a}{1}{2}\) [> length=2],
                },
                ta l3s1 / [coordinate] -- {
                    ta l2s1 / [coordinate] -- {
                        ta l1s3 / \(\explar{a}{3}{1}\),
                        ta l1s2 / \(\explar{a}{2}{1}\),
                    },
                    ta l1s1 / \(\explar{a}{1}{1}\) [> length=2],
                },
            },
        };
        \node at (0, -2.3) {\(T_a\)};
    \end{scope}

    %% Arbre T_b
    \begin{scope}[shift={(8, 0)}, color2]
        \graph[
            phylogenetic tree layout,
            sibling distance=.2cm,
            evolutionary unit length=0.44cm,
            nodes={font=\small},
        ]{
            tb l4s1 / [coordinate] -- {
                tb l3s2 / [coordinate] -- {
                    tb l2s2 / [coordinate] -- {
                        tb l1s6 / \(\explar{b}{3}{2}\),
                        tb l1s5 / \(\explar{b}{2}{1}\),
                    },
                    tb l1s4 / \(\explar{b}{1}{3}\) [> length=2],
                },
                tb l3s1 / [coordinate] -- {
                    tb l1s1 / \(\explar{b}{3}{1}\) [> length=2],
                    tb l2s1 / [coordinate] -- {
                        tb l1s3 / \(\explar{b}{1}{2}\),
                        tb l1s2 / \(\explar{b}{1}{1}\),
                    },
                },
            },
        };
        \node at (0, -2.3) {\(T_b\)};
    \end{scope}
\end{tikzpicture}

\vspace{1em}
\begin{tikzpicture}
    %% Arbre T_c
    \begin{scope}[color3]
        \graph[
            phylogenetic tree layout,
            sibling distance=.2cm,
            nodes={font=\small},
        ]{
            tc l3s1 / [coordinate] -- {
                tc l2s1 / [coordinate] -- {
                    tc l1s3 / \(\explar{c}{3}{1}\),
                    tc l1s2 / \(\explar{c}{2}{1}\),
                },
                tc l1s1 / \(\explar{c}{1}{1}\) [> length=2],
            },
        };
        \coordinate (labels) at (0, -2.3);
        \node at (labels-|tc l1s2) {\(T_c\)};
    \end{scope}

    %% Associations
    \node[anchor=north] at (4.25, 0) {\small\begin{tabular}{@{}l@{}l@{}}
        \(A = \{\)
        & \(\{{\color{color1}\explar{a}{1}{1}}, {\color{color2}\explar{b}{1}{1}}\},
            \{{\color{color2}\explar{b}{1}{1}}, {\color{color2}\explar{b}{1}{2}}\},
            \{{\color{color1}\explar{a}{1}{2}}, {\color{color2}\explar{b}{1}{3}}\},\) \\[3pt]
        & \(\{{\color{color1}\explar{a}{2}{2}}, {\color{color2}\explar{b}{2}{1}}\},
            \{{\color{color2}\explar{b}{2}{1}}, {\color{color3}\explar{c}{2}{1}}\},
            \{{\color{color1}\explar{a}{3}{1}}, {\color{color2}\explar{b}{3}{1}}\},\) \\[3pt]
        & \(\{{\color{color1}\explar{a}{3}{2}}, {\color{color2}\explar{b}{3}{2}}\},
            \{{\color{color2}\explar{b}{3}{2}}, {\color{color3}\explar{c}{3}{1}}\}\}\)
    \end{tabular}};

    %% Arbre S
    \begin{scope}[shift={(8, 0)}]
        \graph[
            phylogenetic tree layout,
            sibling distance=.2cm,
            nodes={font=\small\bfseries},
        ]{
            s l3s1 / [coordinate, label={[host]right:α}] -- {
                s l2s1 / [coordinate, label={[host]right:β}] -- {
                    s l1s3 / 3,
                    s l1s2 / 2,
                },
                s l1s1 / 1 [> length=2],
            },
        };
        \coordinate (labels) at (0, -2.3);
        \node at (labels-|s l1s2) {\(S\)};
    \end{scope}
\end{tikzpicture}

\vspace{1.5em}
\begin{tikzpicture}[x=.8cm, y=1cm]
    \draw[separator rule] (0, 3.95) -- ++(14, 0);

    %% Arbre T_a
    \begin{scope}[color1]
        \graph[no placement, nodes={font=\small}]{
            ta l4s1 / [graph node, x=5.5, y=3, label={[host, black, yshift=2pt]left:α}] -- {
                ta l3s1 / [graph node, x=2, y=2, label={[host, black]left:α}] -- {
                    ta l1s1 / \(\explar{a}{1}{1}\) [x=0],
                    ta l2s1 / [graph node, x=3.33, y=1, label={[host, black, yshift=-2pt]left:β}] -- {
                        ta l1s2 / \(\explar{a}{2}{1}\) [x=3],
                        ta l1s3 / \(\explar{a}{3}{1}\) [x=4],
                    },
                },
                ta l3s2 / [graph node, x=9, y=2] -- {
                    ta l1s4 / \(\explar{a}{1}{2}\) [x=6],
                    ta l2s2 / [graph node, x=10.5, y=1] -- {
                        ta l1s5 / \(\explar{a}{2}{2}\) [x=9],
                        ta l1s6 / \(\explar{a}{3}{2}\) [x=12],
                    },
                },
            },
        };
    \end{scope}

    %% Arbre T_b
    \begin{scope}[color2]
        \graph[no placement, nodes={font=\small}]{
            tb l4s1 / [graph node, x=6.5, y=3] -- {
                tb l3s1 / [graph node, x=3, y=2] -- {
                    tb l1s1 / \(\explar{b}{3}{1}\) [x=5],
                    tb l2s1 / [graph node, x=1.66, y=1, label={[host, black]right:1}] -- {
                        tb l1s2 / \(\explar{b}{1}{1}\) [x=1],
                        tb l1s3 / \(\explar{b}{1}{2}\) [x=2],
                    },
                },
                tb l3s2 / [graph node, x=10, y=2, label={[host, black]right:α}] -- {
                    tb l1s4 / \(\explar{b}{1}{3}\) [x=7],
                    tb l2s2 / [graph node, x=11.5, y=1] -- {
                        tb l1s5 / \(\explar{b}{2}{1}\) [x=10],
                        tb l1s6 / \(\explar{b}{3}{2}\) [x=13],
                    },
                },
            },
        };
    \end{scope}

    %% Arbre T_c
    \begin{scope}[color3]
        \graph[no placement, nodes={font=\small}]{
            tc l3s1 / [graph node, x=11, y=2, label={[host, black]right:α}] -- {
                tc l1s1 / \(\explar{c}{1}{1}\) [x=8],
                tc l2s1 / [graph node, x=12.5, y=1, label={[host, black, yshift=2pt]right:β}] -- {
                    tc l1s2 / \(\explar{c}{2}{1}\) [x=11],
                    tc l1s3 / \(\explar{c}{3}{1}\) [x=14],
                },
            },
        };
    \end{scope}

    %% Adjacences
    \begin{pgfonlayer}{hostedges}
        \path[adjacency] (ta l4s1) edge (tb l4s1);
        \path[adjacency] (ta l3s1) edge (tb l3s1);
        \path[adjacency] (ta l3s2) edge (tb l3s2);
        \path[adjacency] (ta l2s2) edge (tb l2s2);
        \path[adjacency] (tb l2s2) edge (tc l2s1);

        \path[adjacency] (ta l1s1) edge (tb l1s2);
        \path[adjacency] (tb l1s2) edge (tb l1s3);
        \path[adjacency] (ta l1s3) edge (tb l1s1);
        \path[adjacency] (ta l1s4) edge (tb l1s4);
        \path[adjacency] (ta l1s5) edge (tb l1s5);
        \path[adjacency] (tb l1s5) edge (tc l1s2);
        \path[adjacency] (ta l1s6) edge (tb l1s6);
        \path[adjacency] (tb l1s6) edge (tc l1s3);
    \end{pgfonlayer}
\end{tikzpicture}

    \caption[Instance of and solution to an adjacency reconciliation problem]{
        \emph{(Top.)}
        Forest of phylogenetic trees \(\gtrees = T_a \sqcup T_b \sqcup T_c\)
        and phylogeny \(S\).
        Each leaf of \(\gtrees\) is annotated with its host as a subscript
        and a copy number as a superscript, which makes it possible to distinguish it in the set of adjacencies \(\lladjset\).
        \emph{(Bottom.)}
        Adjacency reconciliation of \(\gtrees\) with \(S\).
        Each pair of adjacent nodes is connected by a dotted line.
        The reconciliation of the trees with \(S\) is indicated by labels.
    }
    \label{fig:reconciliation:adj-rec}
\end{figure}

A set of ancestral adjacencies implies temporal constraints that must be compatible with those imposed by the topology of \(\gtrees\).
Thus, a node cannot be adjacent to one of its descendants or ancestors and, if two nodes are adjacent, then an ancestor of one cannot be adjacent to a descendant of the other.

\vspace{1ex}
\begin{definition}[Compatibility]\label{definition:adjrec-adj-compat}
    Let \(\gtrees\) be a phylogenetic forest, \(S\) a phylogenetic tree, and \(\adjset\) a set of adjacencies of \(\gtrees\).
    The adjacencies of \(\adjset\) are \emph{compatible} if there exists a preorder~\(\succeq_\alpha\) on the nodes of \(\tnodes(\gtrees)\) such that:
    \begin{condition}
        \item \(v \succ_\alpha v\child\) for each branch \((v, v\child) \in \tedges(\gtrees)\);
        \item \(v \succeq_\alpha w\) and \(w \succeq_\alpha v\) for each adjacency \(\{v, w\} \in \adjset\). \qedhere
    \end{condition}
\end{definition}

Moreover, for them to correspond to a known biological structure, these adjacencies must be able to be assembled into segments or cycles without forming branchings.

\begin{definition}[Linearity]\label{definition:adjrec-lin}
    Let \(\gtrees\) be a phylogenetic forest, \(S\) a phylogenetic tree, and \(\tup{\rec, \adjset}\) an adjacency reconciliation of \(\gtrees\) with \(S\).
    The adjacencies of \(\adjset\) are \emph{linear} if each \(v \in \tnodes(\gtrees)\) appears in at most two adjacencies.
\end{definition}

As for bilateral reconciliations, choosing an evolutionary model makes it possible to interpret an adjacency reconciliation.
Such a model specifies not only the events that individually affect the dependent units, but also those that affect their adjacencies.
An interpretation is represented by an adjacency forest.

\begin{definition}[Adjacency forests]\label{definition:adjrectree}
    Let \(\gtrees\) be a phylogenetic forest, \(S\) a phylogenetic tree, \(\tup{\rec, \adjset}\) an adjacency reconciliation of \(\gtrees\) with \(S\), and \(\tup{\recmodel\rmunit, \recmodel\rmadj}\) two sets of event types.
    Let \(\rectree(\rec)\) be a reconciled forest under the model \(\recmodel\rmunit\) containing the maps \(\exrec\) and \(\evt\).
    An \emph{adjacency forest} \(\adjtree\) is composed of nodes
        labelled by adjacencies
        \(\exadj : \tnodes(\adjtree) \onto \alladjs(\exrec)\)
        and by events
        \(\evt\rmadj : \tnodes(\adjtree) \onto \recmodel\rmadj \cup \{\evtE, \evtB\}\),
        such that:
    \begin{condition}
        \item only the leaves of \(\adjtree\) can be labelled by the events \(\evtE\) or \(\evtB\), where \(\evtE\) labels the extant adjacencies and \(\evtB\) the adjacency breaks;
        \item \(\lladjset = \{\exadj(v) \suchthat v \in \tnodes(\adjtree)
            \text{ and } \evt(v) = \evtE \}\);
        \item \(\adjset \subseteq \{\exadj(v) \suchthat v \in \tnodes(\adjtree)
            \text{ and } \evt(v) \neq \evtB \}\). \qedhere
    \end{condition}
\end{definition}

Each root of an adjacency forest corresponds to the appearance of a new adjacency between two units\footnote{The model of \textcite{Berard2012EvolutionGeneNeighborhoods} makes it possible, under certain conditions, not to count some roots of the forest as new adjacencies but rather as tandem duplications.}, which then propagates according to the events of \(\recmodel\rmadj\).
This adjacency gain event is denoted~\(\evtG\).
Consequently, a forest without any branch corresponds to a scenario in which all adjacencies appeared spontaneously, whereas a forest made up of a single tree corresponds to the emergence of all adjacencies from a single ancestor.
The cost associated with an adjacency forest is the sum of the costs of its events, including as many gain events as it has roots.

\begin{problem}{\(\tup{\recmodel\rmunit, \recmodel\rmadj}\)-parsimony}
    \probinput{Phylogenies \(\gtrees\) and \(S\);
               leaf association \(\llmap : \tleavs(\gtrees) \onto \tleavs(S)\);
               adjacencies between leaves \(\lladjset \subseteq \alladjs(\llmap)\);
               costs \(\cost : \recmodel\rmunit \cup \recmodel\rmadj \cup \{\evtG\} \onto \nat\).}
    \proboutput{Minimum cost of an acyclic reconciled forest of \(\gtrees\) with \(S\) under the model \(\recmodel\rmunit\) and of an associated adjacency forest under the model \(\recmodel\rmadj\) for the costs \(\cost\).}
\end{problem}

A simple model for the evolution of adjacencies, based on duplication and loss events, was proposed by \textcite{Berard2012EvolutionGeneNeighborhoods}.
In this model, \(\recmodelDL\rmunit = \{\evtD, \evtL, \evtS\}\), \(\recmodelDL\rmadj = \{\evtDD, \evtD, \evtLL, \evtSS\}\), and each node \(v \in \tnodes(\adjtree)\) of a reconciled forest is such that:
\begin{compactitemize}
    \item if \(\evt\rmadj(v) = \evtDD\) (resp. \(\evtSS\)), then
        \(\exadj(v) = \{x, y\}\),
        \(\evt(x) = \evt(y) = \evtD\) (resp. \(\evtS\)), and
        \(\{\exadj(v\lchild), \exadj(v\rchild)\} = \{\{x\lchild, y\lchild\}, \{x\rchild, y\rchild\}\}\)
        or \(\{\exadj(v\lchild), \exadj(v\rchild)\} = \{\{x\lchild, y\rchild\}, \{x\rchild, y\lchild\}\}\);
    \item if \(\evt\rmadj(v) = \evtLL\), then \(v\) is a leaf,
        \(\exadj(v) = \{x, y\}\),
        and \(\evt(x) = \evt(y) = \evtL\);
    \item if \(\evt\rmadj(v) = \evtD\), then \(v\) is unary,
        \(\exadj(v) = \{x, y\}\),
        \(\evt(x) = \evtD\),
        and \(\exadj(v\child) = \{x', y\}\),
        where \(x' \in \children(x)\).
\end{compactitemize}

\Cref{fig:reconciliation:adj-forest} gives an example of a reconciled forest and of an adjacency forest interpreting the adjacency reconciliation of \cref{fig:reconciliation:adj-rec} under the model \(\tup{\recmodelDL\rmunit, \recmodelDL\rmadj}\).

\begin{figure}[htb]
    \capstart
    \centering
    %! TeX root = ../main.tex
% vim: set spelllang=fr:
\begin{tikzpicture}[
    label distance=-2pt,
    reference/.style={font=\footnotesize, opacity=.5},
]
    %% Réconciliation T_a
    \begin{scope}[x=.5cm, y=.8cm]
        \graph[no placement, nodes={dependent}]{
            tal4s1 / [duplication, x=3.5, y=2.75, label={[reference]left:\(\bar{a}_1\)}] -- {
                tal3s1 / [speciation, x=3, y=2, label={[reference]left:\(\bar{a}_2\)}] -- {
                    tal2s1 / [graph node, x=1.5, y=1, label={[reference, yshift=1pt]left:\(\bar{a}_4\)}] --
                        tal1s1 / \(\explar{a}{1}{1}\) [x=0, y=0, >shorten >=-3pt],
                    tal2s2 / [speciation, x=5.75, y=1.1, label={[reference]left:\(\bar{a}_5\)}] -- {
                        tal1s2 / \(\explar{a}{2}{1}\) [x=4, y=0, >shorten >=-3pt],
                        tal1s3 / \(\explar{a}{3}{1}\) [x=7.5, y=0, >shorten >=-3pt],
                    },
                },
                tal3s2 / [speciation, x=4, y=2, label={[reference]right:\(\bar{a}_3\)}] -- {
                    tal1s4 / \(\explar{a}{1}{2}\) [x=1, y=0, >shorten >=-3pt],
                    tal2s3 / [speciation, x=6.75, y=1.1, label={[reference]right:\(\bar{a}_6\)}] -- {
                        tal1s5 / \(\explar{a}{2}{2}\) [x=5, y=0, >shorten >=-3pt],
                        tal1s6 / \(\explar{a}{3}{2}\) [x=8.5, y=0, >shorten >=-3pt],
                    },
                },
            },
        };
    \end{scope}

    %% Réconciliation T_b
    \begin{scope}[x=.5cm, y=.8cm, shift={(12, 0)}]
        \graph[no placement, nodes={dependent}]{
            tbl4s1 / [duplication, x=3.5, y=2.75, label={[reference]left:\(\bar{b}_1\)}] -- {
                tbl3s1 / [speciation, x=3, y=2, label={[reference, yshift=1pt]left:\(\bar{b}_2\)}] -- {
                    tbl2s1 / [duplication, x=1.5, y=1, label={[reference, yshift=1pt]left:\(\bar{b}_4\)}] -- {
                        tbl1s1 / \(\explar{b}{1}{1}\) [x=-.5, y=0, >shorten >=-3pt],
                        tbl1s2 / \(\explar{b}{1}{2}\) [x=.5, y=0, >shorten >=-3pt],
                    },
                    tbl2s2 / [speciation, x=5.75, y=1.1, label={[reference, xshift=-1pt]left:\(\bar{b}_5\)}] -- {
                        tbl1s3 / [loss, >loss branch, rotate=-45, x=4, y=0],
                        tbl1s4 / \(\explar{b}{3}{1}\) [x=7.5, y=0, >shorten >=-3pt],
                    },
                },
                tbl3s2 / [speciation, x=4, y=2, label={[reference, yshift=1pt]right:\(\bar{b}_3\)}] -- {
                    tbl1s5 / \(\explar{b}{1}{3}\) [x=1.5, y=0, >shorten >=-3pt],
                    tbl2s3 / [speciation, x=6.75, y=1.1, label={[reference]right:\(\bar{b}_6\)}] -- {
                        tbl1s6 / \(\explar{b}{2}{1}\) [x=5, y=0, >shorten >=-3pt],
                        tbl1s7 / \(\explar{b}{3}{2}\) [x=8.5, y=0, >shorten >=-3pt],
                    },
                },
            },
        };
    \end{scope}

    %% Réconciliation T_c
    \begin{scope}[x=.5cm, y=1.1cm, shift={(24, 0)}]
        \graph[no placement, nodes={dependent}]{
            tcl3s1 / [speciation, label={[reference]right:\(\bar{c}_1\)}, x=1.5, y=2] -- {
                tcl1s1 / \(\explar{c}{1}{1}\) [x=0, y=0, >shorten >=-2pt],
                tcl2s1 / [speciation, label={[reference]right:\(\bar{c}_2\)}, x=3, y=1.1] -- {
                    tcl1s2 / \(\explar{c}{2}{1}\) [x=2, y=0, >shorten >=-2pt],
                    tcl1s3 / \(\explar{c}{3}{1}\) [x=4, y=0, >shorten >=-4pt],
                },
            },
        };
    \end{scope}

    \begin{pgfonlayer}{hostnodes}
        \node[host leaf node=color1, fit=(tal1s1)(tal1s4)] (sal1s1) {};
        \node[host leaf node=color1, fit=(tal1s2)(tal1s5)] (sal1s2) {};
        \node[host leaf node=color1, fit=(tal1s3)(tal1s6)] (sal1s3) {};
        \node[host node=color1, fit=(tal2s2)(tal2s3)] (sal2s1) {};
        \node[host node=color1, fit=(tal3s1)(tal3s2)] (sal3s1) {};

        \node[host leaf node=color2, fit=(tbl1s1)(tbl1s2)(tbl1s5), inner xsep=-.15cm] (sbl1s1) {};
        \node[host leaf node=color2, fit=(tbl1s3)(tbl1s6)] (sbl1s2) {};
        \node[host leaf node=color2, fit=(tbl1s4)(tbl1s7)] (sbl1s3) {};
        \node[host node=color2, fit=(tbl2s2)(tbl2s3)] (sbl2s1) {};
        \node[host node=color2, fit=(tbl3s1)(tbl3s2)] (sbl3s1) {};

        \node[host leaf node=color3, fit=(tcl1s1), inner xsep=.05cm] (scl1s1) {};
        \node[host leaf node=color3, fit=(tcl1s2), inner xsep=.05cm] (scl1s2) {};
        \node[host leaf node=color3, fit=(tcl1s3), inner xsep=.05cm] (scl1s3) {};
        \node[host node=color3, fit=(tcl2s1), inner xsep=.15cm] (scl2s1) {};
        \node[host node=color3, fit=(tcl3s1), inner xsep=.15cm] (scl3s1) {};
    \end{pgfonlayer}

    \begin{pgfonlayer}{hostedges}
        \path[host edge=color1] (sal1s1.-10) -- (sal3s1.-10) -- (sal3s1.170) -- (sal1s1.170);
        \path[host edge=color1] (sal1s2.-10) -- (sal2s1.-10) -- (sal2s1.170) -- (sal1s2.170);
        \path[host edge=color1] (sal2s1.25) -- (sal3s1.10) -- (sal3s1.190) -- (sal2s1.190);
        \path[host edge=color1] (sal1s3.10) -- (sal2s1.10) -- (sal2s1.190) -- (sal1s3.190);
        \path[host edge terminal=color1] (sal3s1.east) -- ($(sal3s1.east)+(0, 1)$)
            -- ($(sal3s1.west)+(0, 1)$) -- (sal3s1.west);

        \path[host edge=color2] (sbl1s1.-10) -- (sbl3s1.-10) -- (sbl3s1.170) -- (sbl1s1.170);
        \path[host edge=color2] (sbl1s2.-10) -- (sbl2s1.-10) -- (sbl2s1.170) -- (sbl1s2.170);
        \path[host edge=color2] (sbl2s1.25) -- (sbl3s1.10) -- (sbl3s1.190) -- (sbl2s1.190);
        \path[host edge=color2] (sbl1s3.10) -- (sbl2s1.10) -- (sbl2s1.190) -- (sbl1s3.190);
        \path[host edge terminal=color2] (sbl3s1.east) -- ($(sbl3s1.east)+(0, 1)$)
            -- ($(sbl3s1.west)+(0, 1)$) -- (sbl3s1.west);

        \path[host edge=color3] (scl1s1.-10) -- (scl3s1.-10) -- (scl3s1.170) -- (scl1s1.170);
        \path[host edge=color3] (scl1s2.-10) -- (scl2s1.-10) -- (scl2s1.170) -- (scl1s2.170);
        \path[host edge=color3] (scl2s1.25) -- (scl3s1.10) -- (scl3s1.190) -- (scl2s1.190);
        \path[host edge=color3] (scl1s3.10) -- (scl2s1.10) -- (scl2s1.190) -- (scl1s3.190);
    \end{pgfonlayer}
\end{tikzpicture}

%% Forêt d’adjacences
\vspace{1.5em}
\begin{tikzpicture}[reference/.style={font=\small, inner sep=0pt, outer sep=0pt, outer ysep=3pt}]
    \draw[separator rule] (-3.5, .75) -- ++(14, 0);

    \begin{scope}
        \graph[
            phylogenetic tree layout,
            evolutionary unit length=0.55cm,
            nodes={reference},
        ] {
            / {(\(\bar{a}_1\), \(\bar{b}_1\))} -- {
                / {(\(\bar{a}_3\), \(\bar{b}_3\))} -- {
                    / {(\(\bar{a}_6\), \(\bar{b}_6\))} -- {
                        / {\((\explar{a}{3}{2}, \explar{b}{3}{2})\)},
                        / {\((\explar{a}{2}{2}, \explar{b}{2}{1})\)},
                    },
                    / {\((\explar{a}{1}{2}, \explar{b}{1}{3})\)} [> length=2],
                },
                / {(\(\bar{a}_2\), \(\bar{b}_2\))} -- {
                    / {(\(\bar{a}_5\), \(\bar{b}_5\))} -- {
                        / {\((\explar{a}{3}{1}, \explar{b}{3}{1})\)},
                        / \(\strut\star\),
                    },
                    / {(\(\bar{a}_4\), \(\bar{b}_4\))} -- / {\((\explar{a}{1}{1}, \explar{b}{1}{1})\)},
                },
            },
        };
    \end{scope}

    \begin{scope}[shift={(6, -1.87)}]
        \node[reference] {\((\explar{b}{1}{1}, \explar{b}{1}{2})\)};
    \end{scope}

    \begin{scope}[shift={(9, -.57)}]
        \graph[
            phylogenetic tree layout,
            evolutionary unit length=1.05cm,
            nodes={reference},
        ] {
            / {(\(\bar{b}_6\), \(\bar{c}_2\))} -- {
                / {\((\explar{b}{3}{2}, \explar{c}{3}{1})\)},
                / {\((\explar{b}{2}{1}, \explar{c}{2}{1})\)},
            },
        };
    \end{scope}
\end{tikzpicture}

    \caption[Interpretation of an adjacency reconciliation]{
        \emph{(Top.)}
        Reconciled trees of \(T_a\), \(T_b\) and \(T_c\) with \(S\)
        under the \(\recmodelDL\) model
        and according to the reconciliation represented in \cref{fig:reconciliation:adj-rec}.
        The grey labels make it possible to identify the internal nodes in the lower part of the figure.
        \emph{(Bottom.)}
        Adjacency forest representing a \(\tup{\recmodelDL\rmunit, \recmodelDL\rmadj}\) evolutionary scenario starting from three ancestral adjacencies.
        The star \(\star\) corresponds to a break of the adjacency through the loss of one of its endpoints.
    }
    \label{fig:reconciliation:adj-forest}
\end{figure}

Provided that all events except \(\evtG\) and \(\evtB\) are given a zero cost, \textcite{Berard2012EvolutionGeneNeighborhoods} showed that the \(\tup{\recmodelDL\rmunit, \recmodelDL\rmadj}\)-parsimony problem can be solved in polynomial time by precomputing individual \(\recmodelDL\) reconciliations and then running a dynamic programming algorithm.
No polynomial algorithm is known to date for the general problem and \textcite{Boussau2020ReconcilingGeneTrees} conjectured that the problem is NP-complete.

Following the initial publication of \textcite{Berard2012EvolutionGeneNeighborhoods}, the model was extended by \textcite{Patterson2013LateralGeneTransfer} and then by \textcite{Duchemin2017DecostarReconstructingAncestral} under the names DeCoLT and DeCoSTAR to take into account the \(\evtT\) and \(\evtTX\) events.

Since adjacencies are reconstructed separately from one another, these algorithms can return a non-linear set of adjacencies (see \cref{definition:adjrec-lin}).
A linearization step can be carried out as post-processing, but the linearization problem is NP-complete in general~\parencite{Manuch2012LinearizationAncestralMultichromosomal}.
The non-linearity of an optimal result can, however, help detect inconsistencies in the input data~\parencite{Tannier2020AncestralGenomeOrganization}.

The DeCo approach assumes that adjacencies evolve independently of one another, which excludes from the modelling segmental events affecting several adjacencies simultaneously.
This characteristic also means that the loss of a unit inside a segment of adjacencies causes this segment to split in two, and requires the gain of a new adjacency between the two segments.
A way to remedy this last limitation was, however, recently presented by~\textcite{Tannier2020AncestralGenomeOrganization}.

\subsection{Segmental reconciliation}\label{sec:reconciliation:multiple:segmental}

To allow the modelling of segmental events simultaneously affecting more than two units, \textcite{Delabre2020EvolutionSegmentalDuplications} introduced, under the name \emph{super-reconciliation}, a method generalizing the notion of adjacency to sets of units of arbitrary size called \emph{syntenies}.
The internal order of these syntenies can be taken into account (segments) or ignored (groups).

\begin{definition}[Segments]\label{definition:seg}
    Let \(L\) and \(M\) be two sets and \(f : L \onto M\) a map.
    A \emph{segment} of size~\(k\) is a tuple \(x = (x_1, …, x_k) \in L^k\) such that, for all \(1 \leq i, j \leq k\) with \(i \neq j\), \(x_i \neq x_j\) and \(f(x_i) = f(x_j)\).
    The set of possible segments of arbitrary size is denoted \(\allsegs(f)\).
\end{definition}

\begin{definition}[Groups]\label{definition:cluster}
    Let \(L\) and \(M\) be two sets and \(f : L \onto M\) a map.
    A \emph{group} is a subset \(c = \{c_1, …, c_k\} \subseteq L\) such that, for all \(1 \leq i, j \leq k\), \(f(c_i) = f(c_j)\).
    The set of possible groups is denoted \(\allclstrs(f)\).
\end{definition}

The definitions detailed in the previous section regarding adjacency reconciliations and forests carry over directly to the context of super-reconciliation, to define synteny reconciliations and forests.

This extension makes it possible to specify evolutionary models comprising segmental events, which simultaneously affect sets of units.
The evolutionary model studied in the article builds on \(\recmodelDL\).
During a speciation event~(\(\evtS\)), each member unit of the synteny simultaneously undergoes a codivergence and ends up in the syntenies of both children.
A duplication event~(\(\evtD\)) includes the possibility that only a sub-segment of the initial synteny is copied.
Finally, a loss event~(\(\evtL\)) affects a sub-segment of the initial synteny, forming a new synteny from the concatenation of the two remaining parts.
Formally, for a node \(v \in \tnodes(\segtree)\) of the synteny forest such that \(\exseg(v) = (x_1, …, x_k)\):

\begin{compactitemize}
    \item if \(\evt(\segtree) = \evtS\), then
        \(v\) and all \(x_i\) are binary internal nodes,
        \(\exseg(v\lchild) = (y_1, …, y_k)\) and \(\exseg(v\rchild) = (z_1, …, z_k)\),
        and for all \(1 \leq i \leq k\),
        \(\evt(x_i) = \evtS\)
        and \(\{y_i, z_i\} = \{x_i{}\lchild, x_i{}\rchild\}\);
    \item if \(\evt(\segtree) = \evtD\), then
        \(v\) and all \(x_i\) are also binary,
        and there exist two indices \(1 \leq s \leq t \leq k\) such that
        \(\exseg(v\lchild) = (y_1, …, y_k)\) and \(\exseg(v\rchild) = (z_s, …, z_t)\)
        (or conversely for \(v\lchild\) and \(v\rchild\)),
        and for all \(1 \leq i \leq k\),
        \(\evt(x_i) = \evtD\),
        for \(1 \leq i < s\) and \(t < i \leq k\),
        \(y_i \in \{x_i{}\lchild, x_i{}\rchild\}\),
        and for \(s \leq i \leq t\),
        and \(\{y_i, z_i\} = \{x_i{}\lchild, x_i{}\rchild\}\);
    \item if \(\evt(\segtree) = \evtL\), then
        \(v\) and all \(x_i\) are unary and
        there exist two indices \(1 \leq s \leq t \leq k\) such that
        \(\exseg(v\child) = (y_1, …, y_{s-1}, y_{t+1}, …, y_k)\),
        for all \(1 \leq i < s\) and \(t < i \leq k\),
        \(\evt(x_i) = \evtP\) and
        \(y_i = x_i{}\child\),
        and for \(s \leq i \leq t\),
        \(\evt(x_i) = \evtL\).
\end{compactitemize}

The combinatorial optimization problem associated with super-reconciliation is similar to the \(\tup{\recmodel\rmunit, \recmodel\rmadj}\)-parsimony problem defined previously, except that the expected output is a forest of syntenies instead of adjacencies.
This problem comes in a \enquote{large} version, in which the topology of the forest and the content of its syntenies must be inferred simultaneously, and a \enquote{small} version in which the topology is fixed.

In the article, a restriction of the optimization problem is studied in which the synteny forest is made up of a single tree and all segments derive from a unique ancestral segment.
In its large version, the problem is shown to be NP-complete.
The complexity status of the small version remains unknown in the case of segments, whereas the problem is polynomial in the case of groups.
